\documentclass[12pt]{article}
\usepackage[utf8]{inputenc}
\usepackage[T1]{fontenc}
\usepackage{hyperref}
\title{Pallanganmiddang Language Resources for K-12 Curriculum Development}
\author{Compiled by Claude Sonnet 5.5 (claude-sonnet-5-5) for Andr\'as Kornai}
\date{}
\begin{document}
\maketitle

\section{Introduction}
Pallanganmiddang (ISO 639-3 pmd) is an extinct, poorly attested language of north-east Victoria, Australia. According to the Victorian Aboriginal Corporation for Languages (VACL), it is known from four sources (Robinson, Smyth 1878, Curr 1886--7, Murdock 1900) that together give a vocabulary of over 300 words and almost no grammar. Some scholars treat the name as a local group (clan) within the Waywurru language area rather than as a separate language.

\section{Primary Educational Resources}
No school materials were found.

\section{Linguistic and Reference Materials}
\begin{itemize}
\item Wikipedia overview: \url{https://en.wikipedia.org/wiki/Pallanganmiddang_language}
\item Blake and Reid's analysis of the sources, cited as the main reference (Aboriginal History vol.\ 23, 1999): \url{https://press-files.anu.edu.au/downloads/press/p72831/pdf/article0217.pdf}
\item Dhudhuroa and Yaithmathang languages and social groups in north-east Victoria, a reconstruction (ANU Press): \url{https://press-files.anu.edu.au/downloads/press/p74631/html/ch09.html?page=12}
\item ASJP wordlist for Pallanganmiddang: \url{https://asjp.clld.org/languages/PALLANGANMIDDANG}
\end{itemize}

\section{Religious and Community Resources}
\begin{itemize}
\item VACL page on the Waywurru language, the area within which Pallanganmiddang is sometimes placed: \url{https://vacl.org.au/language/waywurru/}
\item Wikipedia page on the Pallanganmiddang people: \url{https://en.wikipedia.org/wiki/Pallanganmiddang_people}
\end{itemize}

\section{Digital and Audio Resources}
No audio or software resources were found.

\section{Curriculum Development Context}
The Victorian Aboriginal languages curriculum document \url{https://vcaa.vic.edu.au:443/Documents/alcv/History.pdf} appeared in the results and may help frame a heritage unit; I did not read it in detail. With only about 300 recorded words, the realistic scope is a revitalization or heritage unit built on the wordlist, developed with the relevant Aboriginal community and VACL, and not a full language-medium curriculum.
\end{document}